\section{Appendix E: Response to Supervisor Comments}\label{sec:appendix-response-to-comments}

This appendix reproduces each comment raised by the supervisor (\textbf{AK}) on the previous draft, together with its resolution (\textbf{GZ}). Each comment was treated as a required change rather than a point to be argued: where a comment identified something as wrong or out of scope, the material was corrected or removed, and where a comment could not be acted on within the uncoded, single-relay scope of this work, that is stated plainly and the item is carried into future work rather than answered with a justification. Because the revision restructured the thesis around a single canonical setup (Section~\ref{sec:introduction}) and, in doing so, \emph{removed} several sections that the earlier draft contained, resolutions that entailed deletion are flagged explicitly with \textbf{[SECTION REMOVED]} so that the change is transparent.

\subsection*{Chapter 1: Introduction}

\begin{enumerate}

\item \AK{Abstract is missing.}\\
\GZ{Resolved. The English and Hebrew abstracts are included via \texttt{chapters/frontmatter.tex} and \texttt{chapters/hebrew\_abstract.tex} respectively, and both were rewritten to match the final (canonical) scope.}

\item \AK{There are compilation errors. Would be good to resolve them.}\\
\GZ{Resolved. The missing-brace errors and a duplicate \texttt{\textbackslash listoflistequations} were fixed; the project compiles cleanly under XeLaTeX with zero undefined references.}

\item \AK{There are compilation errors. Would be good to resolve them. (second instance)}\\
\GZ{Resolved together with the item above; the second flagged location shared the same root cause.}

\item \AK{Unusual citation.}\\
\GZ{Resolved. The inline \texttt{source: [\ldots]} annotations placed under equations were removed throughout, and equation attributions now use standard \texttt{\textbackslash cite} references (e.g.\ Proakis, Goodfellow et al.).}

\item \AK{This is not causal\ldots{} [regarding the symmetric relay window $x_R[i]=f_\theta(y_R[i-w:i+w])$]}\\
\GZ{Accepted; the comment is correct and the earlier text did not say so. The window is \emph{not} causal, and the thesis no longer presents it as though it were. Remark~\ref{rem:window-causality} now labels it non-causal outright, states that it is admissible here only because the relay is half-duplex and buffers the whole listening-phase block before transmitting, quantifies the cost as a fixed $w$-symbol processing latency, and records the non-causality as a limitation of the learned relays rather than a property to be justified. The strictly causal alternative, a past-only window $y_R[i-2w:i]$, is identified as the change any streaming deployment would require; it was not evaluated here, and no claim in the thesis is made for it. AF and DF are unaffected, using $w=0$.}

\item \AK{This is the transmitted signal by the relay.}\\
\GZ{Resolved; the surrounding sentence was corrected and the time index restored.}

\item \AK{The time axis is missing. Aggregation of information can be useful, especially when incorporating an error correction code. Do you allow only scalar operations at the relay?}\\
\GZ{Resolved. The time index is restored, and the relay is explicitly \emph{not} restricted to scalar processing: learned relays use a length-$2w+1$ sliding window ($w>0$), while AF and symbol-wise DF are memoryless ($w=0$). The uncoded setting is stated up front; the value of temporal aggregation is in fact a central theme of the principal contribution (Chapter~\ref{sec:unknown-channels}), where the window is what allows the learned relay to equalise unknown intersymbol interference.}

\item \AK{I agree that AF is memoryless but DF is not. Information-theoretic DF works in blocks to incorporate a strong (capacity-achieving over a single link) code and recovers the message at the end of the time block. You may of course use simple slicing, but this won't do much good.}\\
\GZ{Addressed and formalised as Remark~\ref{rem:df-terminology}. The thesis studies \emph{symbol-wise} (slicing) DF and now states this explicitly wherever DF appears, distinguishing it from information-theoretic block-DF. The distinction is load-bearing in the results: the high-SNR optimality of symbol-wise DF (H2) is explicitly \emph{not} claimed to carry over to the coded block-DF setting, and a coded block-DF comparison is identified as the leading item of future work (Chapter~\ref{sec:discussion-and-conclusions}).}

\end{enumerate}

\subsection*{Chapter 2: Literature Review}

\begin{enumerate}

\item \AK{You haven't defined the system model so this discussion seems out of place.}\\
\GZ{Resolved. The system-model definition was moved to Chapter~\ref{sec:introduction}; Chapter~\ref{sec:literature-review} is now a pure literature review, and the MIMO background section carries a one-sentence pointer noting it supports only the deferred material.}

\item \AK{The relevance to this equation is not clear. Also, what do you mean by scaling? Isn't that the capacity for any SNR?}\\
\textbf{[SECTION REMOVED]} \GZ{The capacity/scaling passage was part of the MIMO capacity discussion. In the canonical restructure the entire MIMO thread was removed from the main line of the thesis (see comment~4 below), so the unclear equation no longer appears. The residual capacity material retained as background was corrected to avoid the ambiguous ``scaling'' phrasing.}

\item \AK{Here you introduce fading for the first time. You could've done that already in the SISO case. Also, now the channels are complex whereas up until now they were real (and with fixed known gains).}\\
\GZ{Resolved. The model conventions were unified: the canonical channel is i.i.d.\ Rayleigh fast fading in complex baseband from the outset. For SISO Hop~1, BPSK is transmitted on the real axis and the receiver applies coherent compensation and takes $\Re\{\cdot\}$, which is the sufficient statistic; this reconciles the earlier real-valued derivations with the complex model and is stated once in Chapter~\ref{sec:introduction}.}

\item \AK{Why not recovering both then? You keep jumping between multiple settings: no-fading/slow-fading/fast-fading, real/complex, SISO/MIMO. Please define one model that you've tackled and stick to it, unless you've actually tackled several settings. Then, each of them deserves a separate chapter.}\\
\textbf{[SECTION REMOVED]} \GZ{Accepted in full; this comment drove the restructure of the thesis. One setup is now fixed and never departed from in the core: SISO on both hops, i.i.d.\ Rayleigh fast fading, complex baseband, QPSK, uncoded BER, with the relay function as the only variable (Table~\ref{tbl:configurations}). The settings the comment lists as being jumped between are gone rather than reconciled. Fading: only ergodic fast fading remains, with no-fading and slow-fading removed. Real versus complex: the model is complex baseband throughout, and a real constellation is never placed on the complex channel. SISO versus MIMO: the $2\times2$ MIMO second hop and its ZF/LMMSE/SIC equalization are removed from the thesis entirely.

Two further settings are each given their own chapter, as the comment instructs. 16-QAM requires a different relay formulation, per-axis processing having broken down, and is Chapter~\ref{sec:extensions}. Unknown and mismatched channels are Chapter~\ref{sec:unknown-channels}. Each states its own configuration at the outset and neither is mixed into the core comparison.

AWGN is present in one role only: it is the channel with exact closed forms, so it is where the simulator is validated against theory (Section~\ref{sec:channel-model-validation-simulative}), paired with BPSK because BPSK is the constellation those closed forms describe. No relay is compared, ranked or selected on AWGN anywhere in the core chapters, and no conclusion is drawn from it. An earlier revision did report a BPSK relay comparison on AWGN alongside the canonical one; that comparison has since been removed.

Also removed outright and redesignated future work: the Rician robustness study, 16-PSK, the CSI-injection ablation, and the end-to-end autoencoder.}

\item \AK{The diversity--multiplexing tradeoff (DMT) is irrelevant in the ergodic setting. DMT assumes slow fading and discusses performance in the presence of some outage probability.}\\
\textbf{[SECTION REMOVED]} \GZ{The DMT discussion was removed. The canonical setting is ergodic fast fading with an uncoded-BER metric, for which DMT does not apply; only a single background sentence remains, with no DMT-based claims.}

\item \AK{Why is it needed at all?}\\
\textbf{[SECTION REMOVED]} \GZ{The flagged passage belonged to the MIMO-equalization background that supported the now-removed MIMO study. With the MIMO thread deferred to future work, the passage was deleted; the remaining equalization background is confined to a clearly-labelled subsection and kept only to the extent needed to define terms used in the future-work discussion.}

\item \AK{It's linear MMSE (LMMSE) unless you work with a Gaussian codebook, which I suspect you do not.}\\
\GZ{Resolved. All references to ``MMSE'' equalization were corrected to ``LMMSE'' (linear MMSE), since the system uses a BPSK constellation rather than a Gaussian codebook. Note that the MIMO equalization material itself is now future work (comment~4); the terminology was corrected in the retained background nonetheless.}

\item \AK{Again, you can construct a variant that recovers a block code, cancels it out from the measurements, then recovers the next block code, etc.\ [coded SIC]}\\
\GZ{Acknowledged and scoped. As with block-DF (Ch.~1, comment~8), the thesis studies uncoded, symbol-level processing; coded per-layer SIC is outside the uncoded scope and is named explicitly in the future-work discussion. Since the MIMO study itself is now deferred (comment~4), coded SIC is doubly future work.}

\item \AK{I thought the goal was to replace the various classical strategies with a neural net that would learn what's the best strategy.}\\
\GZ{Accepted. The goal as originally stated is not what the evidence supports, and the thesis's claim was changed accordingly rather than defended. The result is that a learned relay does \emph{not} universally replace classical strategies: on a matched, known channel symbol-wise DF is at least as good at zero parameter cost (H2). What the learned relay \emph{does} provide (and this is the principal contribution (Chapter~\ref{sec:unknown-channels})) is family-agnostic mitigation when the channel is unknown or mismatched: it restores reliable relaying where memoryless classical relays fail, approaches the model-aware optimum (Viterbi) without knowing the channel family, and does so at constant, parallelizable complexity. The framing is therefore ``characterise when learning helps,'' not ``replace classical processing everywhere'', a change made directly in response to this comment.}

\end{enumerate}

\subsection*{Second round: comments on the revised draft}

The comments below were raised inline on the revised draft, in the abstract, Chapter~\ref{sec:introduction}, Chapter~\ref{sec:literature-review} and the front matter. They were previously answered as coloured margin notes in the document itself, which meant the review exchange was printed in the submitted text. Each is now resolved in place and recorded here instead; no annotation remains anywhere outside this appendix.

\begin{enumerate}

\item \AK{There are some compilation errors pertaining to Hebrew.}\\
\GZ{Resolved. The Hebrew abstract compiles under XeLaTeX with \texttt{bidi} and renders right-to-left; the document builds with zero errors and zero undefined references.}

\item \AK{What's a \underline{fixed} Rayleigh fading?}\\
\GZ{Resolved, and the ambiguity behind the question is now addressed rather than just the phrase. ``Fixed Rayleigh fading'' is gone; the abstract says i.i.d.\ Rayleigh \emph{fast} fading applied identically on both hops. The word ``fixed'' was doing two jobs in the draft, and the collision is what made the question necessary: it meant \emph{held constant across experiments} in most places, but beside a channel it invites the reading \emph{static} or \emph{quasi-static}, which is the opposite of what this channel does. Every use that referred to the channel or the setup is reworded to say ``the same in every experiment'' or ``held unchanged'', and Chapter~\ref{sec:introduction} now states the distinction outright: the channel \emph{model} is identical in every experiment, while the coefficients $h_1[i]$ and $h_2[i]$ are redrawn independently for every symbol. No link in this thesis is static or quasi-static. ``Fixed'' is retained only where it is literally correct, as for the three-tap ISI filter of Chapter~\ref{sec:unknown-channels}, whose taps genuinely do not change.}

\item \AK{DF is not symbolwise, unless you simple slicing.}\\
\GZ{Resolved. DF is named \emph{symbol-wise} decode-and-forward at every occurrence in the abstract, and Remark~\ref{rem:df-terminology} separates it from information-theoretic block-DF.}

\item \AK{What's compact?}\\
\GZ{Resolved. The vague word is replaced by the number: ``an MLP of 169 trainable parameters''.}

\item \AK{From a capacity standpoint, DF is optimal. At short blocklengths, DF should be suboptimal.}\\
\GZ{Accepted. The abstract no longer implies DF is optimal in any general sense. It states that the information-theoretic results concern coded transmission at asymptotic blocklengths, that they do not transfer to the uncoded finite-length setting measured here, and that whether a coded or soft-information relay beats DF remains open.}

\item \AK{In the matched case, according to your claims here, DF is optimal. For mismatched setting, why would one use DF?}\\
\GZ{Accepted; the earlier answer defended keeping DF, which was the wrong response. One would \emph{not} use DF under mismatch, and the thesis does not propose it. DF appears in Chapter~\ref{sec:unknown-channels} only as the incumbent being tested, and the result is that it fails: its BER saturates and then rises with transmit power. That failure is the premise of the chapter, not a recommendation.}

\item \AK{Please use AI to polish this abstract.}\\
\GZ{Done. The abstract was rewritten in full.}

\item \AK{Why? DF cuts the code blocklength by half. Can't this be beaten?}\\
\GZ{Accepted. The claim was too strong and is withdrawn. The abstract now scopes DF's advantage to uncoded, per-symbol relaying over the SNR range actually measured, and states explicitly that neither the measurements nor the cited theory establish that DF cannot be beaten by a coded or soft-information relay.}

\item \AK{AF outperforms DF for uncoded symbols for low SNR values. DF is not universally better across all SNR values.}\\
\GZ{Accepted as a general statement, and the abstract no longer claims universality. For the record, no such crossover occurs within this thesis's own measurements: on the canonical channel DF is below AF at every tabulated point, $0.3337$ against $0.4076$ at 0~dB (Table~\ref{tbl:table2}). The mechanism the comment refers to is real, naive always-forward DF propagating its own decision errors at low SNR, and it is why selective DF exists; it simply does not appear in the range evaluated. The abstract is therefore worded as a statement about the results reported here.}

\item \AK{MLSE w.r.t.\ what? Did you assume colored noise or a channel filter? Did you use convolutional codes? The system model in Chapter~1.1 appears memoryless with white noise.}\\
\GZ{Resolved; the comment identified a genuine gap. MLSE was invoked without the channel that makes it meaningful. The abstract now introduces the three-tap intersymbol-interference filter first and then refers to Viterbi MLSE \emph{matched to that filter}. No convolutional code is used and the noise is white throughout; the memory is in the channel filter, not the noise. The canonical system model is indeed memoryless, which is precisely why the ISI study is a separate chapter.}

\item \AK{This is unclear from the abstract. You cannot refer to a model that you've never mentioned.}\\
\GZ{Resolved together with the item above: the ISI model is stated before anything that depends on it.}

\item \AK{Since the channel is complex, you could've sent QPSK instead of BPSK.}\\
\GZ{Accepted, and acted on rather than answered. QPSK \emph{is} now the canonical constellation, precisely because the channel is complex.}

\item \AK{I am aware of that. In the abstract you talked about QPSK and larger QAM constellations, but you do not mention them here. What is the source of this mismatch? If you use complex channels, stick to 2D constellations or avoid talking about them in the sequel as well.}\\
\GZ{Accepted in full. The instruction is followed literally: the complex channel now carries only two-dimensional constellations. QPSK is canonical and 16-QAM is the extension; BPSK, being one-dimensional, is confined to the real AWGN channel used to calibrate the simulator and is never transmitted over the fading channel. Table~\ref{tbl:configurations} states the three pairings. The mismatch the comment identified, between an abstract promising QPSK and a body delivering BPSK, no longer exists.}

\item \AK{Since you assume uncoded transmission, what reason is there to use $\omega > 0$?}\\
\GZ{Accepted; there is none, and the canonical relay no longer uses one. On a memoryless channel $y_R[i]$ is a sufficient statistic, so $w=0$. A window is used only in Chapter~\ref{sec:unknown-channels}, where the channel has memory to span.}

\item \AK{The goal of the ``system model'' is to provide the model. Either drop this part from this section and define explicitly in Chapter 7 or define in detail what is the competing model you compare yourself against.}\\
\GZ{Resolved by taking the first option. The system model no longer justifies a windowed relay by forward-reference; it states the canonical model, in which the relay is memoryless. The windowed relay is defined in Chapter~\ref{sec:unknown-channels} alongside the three-tap channel that motivates it, and the competing model there, pilot-aided Viterbi MLSE over that filter, is specified in full.}

\item \AK{Didn't you assume white noise and i.i.d.\ Rayleigh fading? If so, why do you talk about correlation and how can $w>0$ help?}\\
\GZ{Accepted; the passage was wrong and has been removed. Under white noise and i.i.d.\ fading there is no correlation for a window to exploit, so the claimed variance-reduction benefit did not hold. The window's justification now rests solely on channel \emph{memory}, which exists only in Chapter~\ref{sec:unknown-channels}.}

\item \AK{You need to use square-brackets to display more concise table and figure titles in the corresponding lists.}\\
\GZ{Resolved. Every caption whose full text ran long now carries a short-form optional argument, so the List of Tables and List of Figures show concise titles instead of full paragraphs; the handful of captions that were already a single short phrase are unchanged, since a short form would duplicate them.}

\end{enumerate}

\noindent\textbf{Summary of removed material.} In direct response to comment~4 above, the following were removed from the thesis and redesignated as future work: the Rician channel-robustness study; the $2\times2$ MIMO second hop and its ZF/LMMSE/SIC equalization; 16-PSK modulation; the CSI-injection ablation; and the end-to-end autoencoder. AWGN was \emph{not} removed but reduced to a single role, calibration against closed-form BER, with no relay compared on it; the BPSK relay comparison an earlier revision reported on AWGN was itself removed. Rician, the other channel family in the earlier draft, was removed outright. The diversity--multiplexing tradeoff discussion (comment~5) and the MIMO-capacity/equalization background passages (comments~2 and~6) were likewise deleted. All removals served the single-canonical-model discipline requested in comment~4, and none affects the validity of the retained results, which were re-run and re-derived entirely within the canonical setup.
